% Detailed solutions for the five canonical VI/41 challenges.

% VI41-CHALLENGE-SOLUTION-01
\subsection*{Solution to Challenge VI.41.01 — The localization sequence}
\begin{solution}
Let \(B=S^{-1}A\). Height-one primes of \(B\) correspond exactly to height-one
primes \(\mathfrak p\subset A\) disjoint from \(S\), via
\(\mathfrak p\mapsto S^{-1}\mathfrak p\). Therefore restriction gives a surjection
\[
\rho:\Div(A)\twoheadrightarrow\Div(B)
\]
whose kernel is freely generated by the height-one primes meeting \(S\).

The fraction fields agree. For \(\mathfrak p\cap S=\varnothing\),
\[
B_{S^{-1}\mathfrak p}\cong A_{\mathfrak p},
\]
so their normalized DVR valuations coincide. Hence, for every \(f\in K^\times\),
\[
\rho(\operatorname{div}_A(f))=\operatorname{div}_B(f).
\]
Thus \(\rho\) descends to a surjection \(\Cl(A)\twoheadrightarrow\Cl(B)\).

If a class \([D]\) maps to zero, choose \(f\in K^\times\) with
\(\rho(D)=\operatorname{div}_B(f)\). Then
\[
\rho(D-\operatorname{div}_A(f))=0,
\]
so \(D-\operatorname{div}_A(f)\) is supported only on height-one primes meeting
\(S\). Therefore the kernel is generated by their classes. Consequently
\[
\Cl(S^{-1}A)\cong
\Cl(A)\Big/\left\langle[\mathfrak p]:
\height\mathfrak p=1,\ \mathfrak p\cap S\ne\varnothing\right\rangle.
\]
This proves the full localization exact sequence, including compatibility of
multiplicities in principal divisors.
\end{solution}

% VI41-CHALLENGE-SOLUTION-02
\subsection*{Solution to Challenge VI.41.02 — The quadric cone without parametrization}
\begin{solution}
Let
\[
A=k[x,y,z]/(xy-z^2),\qquad L=V(x,z),\qquad M=V(y,z).
\]
Normality follows from Serre's criterion: the hypersurface is \(S_2\), and it is
regular in codimension one.

At \(L\), \(y\) is a unit and \(x=z^2/y\), so
\[
v_L(z)=1,\qquad v_L(x)=2.
\]
By symmetry,
\[
\operatorname{div}(x)=2L,\qquad
\operatorname{div}(y)=2M,\qquad
\operatorname{div}(z)=L+M.
\]
Thus \(2[L]=0\) and \([L]=[M]\).

Now
\[
A_x\cong k[x,x^{-1},z]
\]
is a UFD, while \(L\) is the unique height-one prime containing \(x\). The
localization theorem therefore shows that \([L]\) generates \(\Cl(A)\).

If \([L]=0\), write \(L=\operatorname{div}(f)\). Then
\[
\operatorname{div}(f^2/x)=0.
\]
The Krull intersection formula implies \(f^2/x\in A^\times\), and the connected
positive grading gives \(A^\times=k^\times\). Thus \(f^2=cx\). But
\(K(A)=k(x,z)\), and the \(x\)-adic valuation of \(cx\) is \(1\), whereas every
square has even valuation. This is impossible. Hence \([L]\ne0\), and
\[
\boxed{\Cl(A)\cong\mathbb Z/2\mathbb Z}.
\]
No parametrization of the cone is used.
\end{solution}

% VI41-CHALLENGE-SOLUTION-03
\subsection*{Solution to Challenge VI.41.03 — Punctured cone and gluing}
\begin{solution}
Put \(U=X\setminus\{o\}\) and \(L=V(x,z)\). On \(D(y)\), the relation
\(x=z^2/y\) gives
\[
(x,z)A_y=(z)A_y,
\]
so \(z\) is a local equation for \(L\). On \(U\setminus L\), use the local
equation \(1\). On the overlap, \(z\) is invertible: if \(y\ne0\) and \(z=0\),
then \(x=0\), placing the point on \(L\), which has been removed. Thus the two
equations differ by a unit, and \(L|_U\) is Cartier.

It is not globally principal. If \(L|_U=\operatorname{div}_U(f)\), then
codimension-two invariance would identify the class of \(L\) with zero in
\(\Cl(X)\), contradicting
\[
[L]\ne0\in\Cl(X)\cong\mathbb Z/2.
\]
Finally \(2L=\operatorname{div}(x)\) restricts to \(U\), so the Cartier class has
order two. VI/42 reinterprets this gluing class as an order-two line bundle.
\end{solution}

% VI41-CHALLENGE-SOLUTION-04
\subsection*{Solution to Challenge VI.41.04 — Projective-space class group}
\begin{solution}
First use \(U=D_+(x_0)\cong\AA^n\) and \(H=V_+(x_0)\). Since
\(\Cl(U)=0\), the open-subset exact sequence gives
\[
\mathbb Z[H]\twoheadrightarrow\Cl(\PP^n).
\]
A principal divisor has degree zero: write its rational function as \(F/G\) with
\(F,G\) homogeneous of equal degree. Thus no nonzero multiple of \(H\) is
principal, and
\[
\Cl(\PP^n)\cong\mathbb Z[H].
\]

For a second proof, use that \(k[x_0,\ldots,x_n]\) is a UFD. Every
codimension-one homogeneous prime is generated by an irreducible homogeneous
polynomial \(F\). If \(\deg F=d\), choose a linear form \(\ell\); then
\[
\operatorname{div}(F/\ell^d)=V_+(F)-dV_+(\ell).
\]
Thus every prime-divisor class equals \(d[H]\). Degree kills principal divisors
and sends \(H\) to \(1\), so degree gives the inverse isomorphism with
\(\mathbb Z\).

Unique factorization enters the second proof in the principal generation and
factorization of codimension-one homogeneous primes; in the first proof it
enters through \(\Cl(\AA^n)=0\).
\end{solution}

% VI41-CHALLENGE-SOLUTION-05
\subsection*{Solution to Challenge VI.41.05 — Local class groups}
\begin{solution}
Let \(X=\Spec A\) with \(A\) normal Noetherian. Localization at a point
\(x\leftrightarrow\mathfrak p_x\) gives
\[
\Cl(A)\longrightarrow\Cl(A_{\mathfrak p_x}).
\]
If a Weil divisor \(D\) is Cartier, then near \(x\) it has one rational equation
\(f_x\); hence its localized divisor is \(\operatorname{div}(f_x)\), so its
local class is zero.

Conversely, suppose the local class of \(D\) vanishes for every \(x\). Then for
each \(x\) there is \(f_x\in K(X)^\times\) whose divisor on
\(\Spec\OO_{X,x}\) equals the localized divisor \(D_x\). Equality of the
finitely many relevant codimension-one coefficients persists after shrinking,
so on some neighborhood \(U_x\),
\[
D|_{U_x}=\operatorname{div}(f_x).
\]
On overlaps, \(f_x/f_y\) has zero order at every codimension-one point.
Normality and the intersection-of-height-one-localizations property imply,
after shrinking, that this quotient is a regular unit. Thus the \(f_x\) are
compatible Cartier local equations, and \(D\) is Cartier.

Hence local class groups measure the local obstruction to a Weil divisor being
Cartier. Once that obstruction vanishes, the remaining question is whether the
local equations glue to one global rational equation; VI/42 packages that
global gluing obstruction as the Picard group.
\end{solution}
